\section{Conclusion}
\label{sec:conclusion}

We asked whether quantum components help exemplar-free class-incremental
learning, and answered with a protocol that resolves a few tenths of a point.
Across seven placements, no quantum component credibly beats the classical
method that receives the same inputs; the one interval above zero is
$+0.04$ points. Three lessons generalise beyond this baseline. First, measure
the noise before claiming a gain: an earlier $+0.36$ of ours vanished, and
evaluating all inference heads on one network cuts the noise of a comparison
to $0.03$--$0.05$ points. Second, test the classical twin: the PGM fusion beat
the baseline but not ridge regression on the same statistics. Third, know which
degrees of freedom are real: circuits after an amplitude encoding and state
shrinkage cannot change the measurement, and repairing a regulariser's measured
weaknesses does not move it past a trade-off that the shared adapter imposes.
The state-discrimination view did produce a useful classical result: the class
statistics kept by exemplar-free methods carry information the aligned
classifier does not use, and a calibrated fusion recovers part of it without
training or new memory.
